Annual 10-K filings help separate more credible AI disclosure from narrative positioning. The classification of AI-related disclosure in 10-K filings is scalable and based on audited reporting content. The main construct, PatentMismatch, combines low-credibility disclosure composition with contemporaneous patent weakness. In the 2016--2025 sample, low-credibility AI disclosure rises sharply after 2022; for a substantial subset of firms, AI language increases more rapidly than observable patenting activity.

Across alternative definitions, the filing-based mismatch measure predicts weaker subsequent AI patent grants. This specification is the most robust version of the construct across tests. In the AI-talking sample, mismatch firm-years subsequently exhibit lower profitability and higher R\&D intensity. The patterns are more consistent with delayed realization and subsequent adjustment in real activity than with purely non-informative disclosure.

Low-credibility disclosure also varies with scrutiny and incentives. Among continuing AI disclosers, disclosure composition becomes more credible following SEC comment-letter scrutiny, while AI focus remains largely unchanged. The SEC's March 18, 2024 AI-washing enforcement actions also coincide with a broader salience result (SEC, 2024). The March 2024 evidence is better read as regulatory salience or acceleration than as a sharply identified enforcement effect. Low-credibility AI disclosure is also more common when CEO equity incentives are stronger. Around large equity-issuance windows, disclosure credibility deteriorates prior to issuance and partly reverts afterward. Low-credibility AI disclosure appears shaped by both monitoring and opportunistic incentives.

The market-pricing results are more limited. Filing-date returns do not cleanly separate mismatch from non-mismatch firms, and post-filing return evidence remains benchmark-sensitive. The strongest return result is a short-horizon equal-weight alpha, with no comparable value-weight evidence. The issue-window evidence suggests a state-dependent rather than uniform short-run return pattern. Market tests therefore support a narrower interpretation than a broad AI-washing anomaly. The stronger evidence is that the filing-based credibility measure predicts later technological realization and organizational behavior, even when market pricing remains incomplete and uneven.

AI-related narratives appear to diffuse faster than real AI implementation, and the gap is widest exactly where it is hardest to check: small firms, financing windows, and the period right after ChatGPT lowered the cost of AI language. What we cannot yet say is whether investors are fooled in a way that moves prices for long. Filing-date returns do not separate the firms, the post-filing evidence is confined to small stocks, and the value-weighted spreads are flat. Whether AI washing is mostly a credibility problem that resolves through later realization, or also a pricing problem that a sharper design would uncover, is the open question we leave for further work.
